\documentclass[11pt]{article}
\usepackage[utf8]{inputenc}
\usepackage[T1]{fontenc}
\usepackage{babel}
\usepackage{lmodern}
\usepackage{amsmath,amssymb}
\usepackage{hyperref}
\usepackage{mathptmx}

% THEOREM CONFIGS (ENGLISH)
\usepackage[thmmarks]{ntheorem}

\theoremstyle{plain}
\theoremheaderfont{\normalfont\bfseries}
\theorembodyfont{\normalfont}
\theoremseparator{:}
\theorempreskip{\topsep}
\theorempostskip{\topsep}

\theoremsymbol{\ensuremath{\bigcirc}}
\theorembodyfont{\normalfont}
\newtheorem{definition}{Definition}[section]

\theoremsymbol{\ensuremath{\triangle}}
\theorembodyfont{\normalfont}
\newtheorem{note}[definition]{Note}

\theoremsymbol{}
\theorembodyfont{\normalfont}
\newtheorem{proposition}[definition]{Proposition}

\theoremsymbol{\ensuremath{\square}}
\theorembodyfont{\normalfont}
\newtheorem{freeprop}[definition]{Proposition}

\theoremstyle{nonumberplain}
\theoremheaderfont{\normalfont\bfseries}
\theorembodyfont{\normalfont}
\theoremsymbol{\ensuremath{\blacksquare}}
\newtheorem{proof}[definition]{Proof}

% Generic Macros (english)
\newcommand{\naturals}{\mathbb{N}}
\newcommand{\integers}{\mathbb{Z}}
\newcommand{\rationals}{\mathbb{Q}}
\newcommand{\reals}{\mathbb{R}}
\newcommand{\complexs}{\mathbb{C}}
\newcommand{\set}[1]{\{#1\}}
\newcommand{\wee}{\wedge}
\newcommand{\U}{\mathbb{U}}
\newcommand{\negset}[1]{\overline{#1}}
\newcommand{\nequiv}{\not\equiv}

\title{Second Order Fun}
\author{Gian Carlo}
\date{\today}

\begin{document}
\maketitle

\section{Introduction}

Some things I have come across nudge me toward second order logic. I wish to see how it is to be working with it, or on it, which is where this text comes from.

\section{Investigation}

The first thing I would look for in a second order logic theory is set theory. For this, we will draw inspiration from ZFC. Specifically, ZFC allows for new sets by sneaking in an axiom scheme that produces a set by first finding the elements in an already existing set, and then evaluating a first order logic formula on these elements. That, together with an observation I was told was by Quine \footnote{that second order logic is set theory in disguise}, has me thinking we will just define membership as a formula on formulas. 

For notation, I will use lower case letters for first order variables and upper case letters for second order variables. The first axiom we will consider is the following: \begin{equation}
    \forall x \forall F(y) : x \in F \iff F(x).
\end{equation}

It is important to note here that $F(y)$ has to mean precisely that $F$ is a formula with only one first-order variable, $y$. From this point onwards, we will call formulas with a single first order variable \textbf{sets}, and the first order variables will be called \textbf{elements}.

The previous axiom defines membership. Note that the only free variable expected in $F$ is a first order variable. We can not substitute $y$ for a formula. This is to avoid Russel's set. In particular, membership is defined only between first order variables as member of second order variables. So a formula like $x \in x$ or $\neg (x \in x)$ is not a valid formula per our definition, since $x$ has to be both a first order and a second order variable. So far, this seems to keep Russel's set at bay.

The way we will obtain the empty set (and a universal set) is by adding to our language two formulas as axioms: $\bot(x)$ and $\top(x)$, given by \begin{equation}
    \forall x : \neg \bot(x)
\end{equation}
and \begin{equation}
    \forall x : \top(x).
\end{equation}
It should be easy to see that $\bot(x)$ yields us the empty set, while $\top(x)$ is a universal set.

Let us define equality (and equivalence) so that we may affirm unicity of both sets. We will define \begin{equation}
    \forall F(x) \forall G(y) : F \equiv G \iff (\forall z : F(z) \iff G(z)).
\end{equation}
In other words, two sets are equivalent if and only if they have the same elements. In synthesis, I will attempt to detail the proofs.

For elements, \begin{equation}
    \forall x \forall y : x = y \iff (\forall F(z) : F(x) \iff F(y)),
\end{equation}
which tells us that two elements are equal (or identical) if and only if we can not distinguish them by any formula.

With this, we should be able to demonstrate that $\bot(x)$ is equivalent to any formula that is always false, and that $\top(x)$ is equivalent to any tautology. As such, we can give unique names to these sets: $\emptyset$ for $\bot(x)$ and $\mathbb{U}$ for $\top(x)$.

From this point on, any second order variable is to be understood as a set. Recall that sets here are second order formulas with only one free first order variable. The usual set operations are very naturally extended from the logical operators, and here the complete formulas are not given but the idea should be clear. \begin{enumerate}
    \item $x \in F \cap G \iff F(x) \wee G(x)$;
    \item $x \in F \cup G \iff F(x) \vee G(x)$;
    \item $x \in F - G \iff F(x) \wee \neg G(x)$;
    \item $x \in \overline{F} \iff \neg F(x)$; and
    \item $F \subseteq G \iff F(x) \implies G(x)$.
\end{enumerate}

From this, we should get that $\negset{\emptyset} = \mathbb{U}$, $\negset{\U} = \emptyset$, $\emptyset \subseteq \U$.

Let us just get some useful definitions out of the way: \begin{enumerate}
    \item $\forall x \forall y : x \neq y \iff \neg (x = y)$; and
    \item $\forall F \forall G : F \nequiv G \iff \neg (F \equiv G)$.
\end{enumerate}

Unless we affirm the existence of at least one element $x \in \U$, we can not show that $U \neq \emptyset$. So let us assume the existence of exactly one element. Since we know it is one, we can name it $c$. Formally, \begin{equation}
    \exists x \forall y : x = y.
\end{equation}

Since we are positing it's existence, we can replace it by $c$ as a constant and obtain a set: \begin{equation}
    \forall y : y = c
\end{equation}
which is clearly satisfied by $c$ and as such $c \in (\forall y : y = c)$. In order to allow for abbreviation, I will allow $\equiv$ to also work for definitions. As such, we can define \begin{equation}
    R(x) \equiv \forall x : x = c
\end{equation}
and then we may formulate $c \in R$.

Here perhaps a discussion ought to be had on how to interpret $F(x)$ for a set $F$. It is meant to be natural: $F(x)$ means $x$ is a first order variable quantified some way, either with $\forall$ or with $\exists$. To know if $F(y)$ for some $y \in \U$, we drop the quantifier and replace $x$ with $y$ in what is left of the formula. So $\exists x : G(x)$ turns into $G(y)$. Since $G(y)$ is just a formula, we should be able to evaluate it.

We can postulate a finite amount of constants in a finite formula. As an example, here is the formula for two constans: \begin{equation}
    \exists x \exists y \forall z : z = x \vee z = y \wee y \neq x.
\end{equation}

This formula grows combionatorially in size the more distinct constants we wish to postulate, but it is always finite.

A first little result we can see is that if we posit exactly $n$ constants, we expect exactly $2^n$ distinct possible sets. Intuitively, these would be the subsets of $\U$.

\section{Synthesis}

To be done.

\section{Not Conclusions}

I think I like second order logic. Sure, we don't have an effective proof method for every truth, but who needs every truth? Do we have an effective observation method for everything in our material universe? Try to observe something smaller than a planck length in 2026.

As far as Löwenheim-Skolem not working, does not seem like much of an issue. How many models do we really care for anyway? 

For compactness, I don't know about you but I do not have the time to write down an infinite amount of axioms that I can not model. In fact I would hope to keep the amount of axioms very finite and rememberable.

I do know there is a lot of reinventig the wheel here. I like reinventing the wheel. I feel it makes me appreciate the wheel.

\end{document}